\chapterrule
\section{矩阵乘法到底在算什么？}
\label{intro:03-matrix-multiplication}

上一章为需求与一篇文章算出一个匹配分数。如果有一万篇候选，
能不能把这些点积一次排好？先沿用原来的三篇文章，看看怎样从一个分数得到一整行分数。

\subsection{把三篇文章排成一张表}
\phantomsection\label{concept:matrix}

仍按“技术、生活”的顺序，每篇文章占一行：
\[
K=\begin{bmatrix}0.8&0.6\\1&0\\0&1\end{bmatrix}.
\]
第一行是文章 A，第二行是 B，第三行是 C。
这张有行和列的数字表叫作\textbf{矩阵（matrix）}。
$K$ 有 3 行、2 列，形状记为 $3\times2$；记形状时\textbf{先数行，再数列}。
这里每行对应一篇文章，每列对应一种特征，所有数字都来自上一章的人工特征表。

\subsection{先把候选变成列：转置有什么用？}
\label{sec:layout-matrix-card}
\phantomsection\label{concept:transpose}

需求仍横着写成 $\vect q^\top=(0.6,0.8)$。
要让它与文章 A 逐项配对，可以把 A 的两个数竖着排。
把整张候选表的行与列交换，就得到
\[
K^\top=\begin{bmatrix}0.8&1&0\\0.6&0&1\end{bmatrix}.
\]
这个操作叫\textbf{转置（transpose）}，用右上角的 $\top$ 表示。
转置后，A、B、C 分别占一列，形状从 $3\times2$ 变成 $2\times3$；数字及其配对关系没有改变。
本书也用 $\vect q^\top$ 表示横写的行向量，$\vect q$ 表示对应的列向量。
上一章两行向量的点积 $\vect q^\top\cdot\vect a^\top$，
也可以写成一行乘一列的 $\vect q^\top\vect a$；两种写法都得到同一个数。

% The Figure structure intentionally includes the non-floating caption.
% The box must stay unbreakable: inside this minipage a breakable tcolorbox
% still splits against \pagegoal-\pagetotal, and the stacked parts overlap.
\par\noindent\begin{minipage}{\linewidth}
\tagstructbegin{tag=Figure,alttext={三条二维向量按行放置形成三乘二矩阵，转置后按列排列得到二乘三矩阵}}
\begin{tcolorbox}[
  enhanced,breakable=false,colback=paper,colframe=linegray,
  boxrule=.7pt,arc=2mm,left=3mm,right=3mm,top=3mm,bottom=3mm,
  fonttitle=\sffamily\bfseries,coltitle=ink,
  title={图解：同一批向量，可以按行排，也可以按列排}
]
\small
这里先不谈 Attention。假设有三条二维向量：
$(1,0)$、$(0,1)$、$(1,1)$。
\textbf{每种底色始终表示同一条向量}，转置只改变行列方向，不改变数字。

\vspace{1.5mm}
\noindent
\begin{minipage}[t]{.49\linewidth}
\centering
\textbf{按行排：三行、两列}\\[2mm]
\begin{tikzpicture}[x=10mm,y=9mm]
\foreach \r/\shade in {1/paper,2/linegray!25,3/linegray!15}{
  \foreach \c in {1,2}{
    \filldraw[draw=linegray,fill=\shade,line width=.55pt]
      (\c-1,1-\r) rectangle (\c,-\r);
  }
}
\node at (.5,-.5){\sffamily 1}; \node at (1.5,-.5){\sffamily 0};
\node at (.5,-1.5){\sffamily 0}; \node at (1.5,-1.5){\sffamily 1};
\node at (.5,-2.5){\sffamily 1}; \node at (1.5,-2.5){\sffamily 1};
\end{tikzpicture}
\\[1mm]
$A:\ 3\times2$\\[-1mm]
{\footnotesize 每行一条向量}
\end{minipage}%
\begin{minipage}[t]{.49\linewidth}
\centering
\textbf{按列排：两行、三列}\\[2mm]
\begin{tikzpicture}[x=10mm,y=9mm]
\foreach \c/\shade in {1/paper,2/linegray!25,3/linegray!15}{
  \foreach \r in {1,2}{
    \filldraw[draw=linegray,fill=\shade,line width=.55pt]
      (\c-1,1-\r) rectangle (\c,-\r);
  }
}
\node at (.5,-.5){\sffamily 1}; \node at (1.5,-.5){\sffamily 0};
\node at (2.5,-.5){\sffamily 1};
\node at (.5,-1.5){\sffamily 0}; \node at (1.5,-1.5){\sffamily 1};
\node at (2.5,-1.5){\sffamily 1};
\end{tikzpicture}
\\[1mm]
$A^\top:\ 2\times3$\\[-1mm]
{\footnotesize 每列一条向量}
\end{minipage}
\[
A^\top=\begin{bmatrix}1&0&1\\0&1&1\end{bmatrix}.
\]

\textbf{先看方向，再检查形状。}
一般地，$M$ 条 $d$ 维向量，按行排列是
$\underbrace{M\times d}_{\text{一条向量占一行}}$，
转置后按列排列是
$\underbrace{d\times M}_{\text{一条向量占一列}}$。
\end{tcolorbox}

\begin{center}\captionof{figure}{行向量与列向量的转置关系}\label{fig:row-column}\end{center}
\tagstructend
\end{minipage}\par


\subsection{一个需求，一次得到三个分数}

让需求行分别与候选的三列做点积，把结果按原顺序排成一行：
\begin{equation}
\begin{aligned}
\vect q^\top K^\top
&=(0.6,0.8)\begin{bmatrix}0.8&1&0\\0.6&0&1\end{bmatrix}\\
&=(0.6\times0.8+0.8\times0.6,\ 0.6\times1+0.8\times0,\ 0.6\times0+0.8\times1)\\
&=\boxed{(0.96,0.6,0.8)}.
\end{aligned}
\end{equation}
这正是上一章逐篇算出的三个分数，排序仍为 A、C、B。
\textbf{行向量乘矩阵，就是让这一行分别与右边的每一列做点积。}
形状也对应：$(1\times2)(2\times3)\to1\times3$。

同一次乘法也可以看成\textbf{按输入系数加权组合矩阵的行}：
\[
0.6(0.8,1,0)+0.8(0.6,0,1)=(0.96,0.6,0.8).
\]
两种看法只是把同一组乘加按不同顺序组织。

\subsection{两个需求，一次得到两行结果}
\label{sec:matrix-article-batch}

现在再来一个稍偏向技术的需求 $(0.8,0.6)$。
把两个需求按行排好，仍与同一批文章比较：
\[
Q=\begin{bmatrix}0.6&0.8\\0.8&0.6\end{bmatrix},\qquad
QK^\top=\boxed{\begin{bmatrix}0.96&0.6&0.8\\1&0.8&0.6\end{bmatrix}}.
\]
第一行回答原需求，第二行回答新需求；每一列始终对应同一篇文章。
例如第二行第一列为 $0.8\times0.8+0.6\times0.6=1$。
\textbf{矩阵乘矩阵，就是把每个输入行都这样算一次。}

% The caption belongs to a semantically tagged Figure.
\tagstructbegin{tag=Figure,alttext={矩阵乘法的行与列配对：左矩阵第一行一二，右矩阵第一列五六，点积十七；左矩阵第二行三四与右第一列点积三十九。}}
\begin{center}
\begin{tikzpicture}[x=12mm,y=9mm,>=Latex]
\node[font=\small\bfseries] at (0,2.0) {左矩阵 $A$};
\node[font=\small\bfseries] at (4.5,2.0) {右矩阵 $B$};
\node[font=\small\bfseries] at (9,2.0) {结果 $AB$};
\draw[rounded corners=1pt,fill=lavender,draw=purple] (-1.5,.65) rectangle (1.6,1.4);
\node at (-.55,1.02) {$1$}; \node at (.65,1.02) {$2$};
\draw[rounded corners=1pt,fill=paper,draw=linegray] (-1.5,-.35) rectangle (1.6,.4);
\node at (-.55,.02) {$3$}; \node at (.65,.02) {$4$};
\draw[rounded corners=1pt,fill=lavender,draw=purple] (3.7,-.4) rectangle (4.7,1.45);
\node at (4.2,1.02) {$5$}; \node at (4.2,.05) {$6$};
\draw[rounded corners=1pt,fill=paper,draw=linegray] (5.0,-.4) rectangle (6.0,1.45);
\node at (5.5,1.02) {$7$}; \node at (5.5,.05) {$8$};
\draw[rounded corners=1pt,fill=lavender,draw=purple] (7.7,.65) rectangle (10.0,1.4);
\node at (8.4,1.02) {$17$};\node at (9.4,1.02) {$23$};
\draw[rounded corners=1pt,fill=paper,draw=linegray] (7.7,-.35) rectangle (10,.4);
\node at (8.4,.02) {$39$}; \node at (9.4,.02) {$53$};
\draw[->,thick,purple] (1.6,1.0) to[out=20,in=160] (7.7,1.02);
\draw[->,thick,muted] (4.7,-.45)--(4.7,-.9)
 node[below,font=\footnotesize] {列 $j=1$};
\node[font=\footnotesize,align=center] at (5.5,-1.9)
 {行 $(1,2)$ $\cdot$ 列 $(5,6)^\top$ $=17$\\
  行 $(3,4)$ $\cdot$ 列 $(5,6)^\top$ $=39$};
\end{tikzpicture}
\end{center}
\begin{center}\captionof{figure}{矩阵乘法逐项定位：输出 $(i,j)$ 是左矩阵第 $i$ 行与右矩阵第 $j$ 列的点积}\label{fig:matrix-row-column-pairing}\end{center}
\tagstructend


\subsection{从例子推广：形状怎样检查？}

一般地，结果第 $i$ 行、第 $j$ 列是左矩阵第 $i$ 行与右矩阵第 $j$ 列的点积：
\[
(AB)_{ij}=\sum_{r=1}^{n}A_{ir}B_{rj},\qquad
\underbrace{A}_{m\times n}\underbrace{B}_{n\times p}
\longrightarrow\underbrace{AB}_{m\times p}.
\]
内侧两个 $n$ 必须相同，点积才能逐项配对；结果形状由外侧的 $m,p$ 决定。
\begin{center}\small
\tagpdfsetup{table/header-rows={1}}
\begin{tabular}{@{}lll@{}}\toprule
要计算的形状 & 能否相乘 & 输出形状\\\midrule
$(1\times2)(2\times3)$ & 可以：一个需求与三篇文章 & $1\times3$\\
$(2\times2)(2\times3)$ & 可以：两个需求与三篇文章 & $2\times3$\\
$(1\times2)(3\times2)$ & 不可以：$2\ne3$ & 无输出\\\bottomrule
\end{tabular}
\end{center}
这叫\textbf{形状检查}。直接写 $\vect q^\top K$ 不合法，正好说明了前面为何需要转置。

\begin{warningbox}
矩阵乘法一般不满足交换律：$AB$ 与 $BA$ 的形状甚至可能不同。
对于形状相容的矩阵，结合律 $(AB)C=A(BC)$ 成立，可以改变计算的分组。
\end{warningbox}

\subsection{从计算匹配分数，到生成新的特征}
\phantomsection\label{concept:linear-map}

上面的右矩阵装的是候选文章。若换成一张可训练的权重矩阵 $W$，
同样的乘法就能从输入特征生成输出特征：
\[
\underbrace{\vect x^\top}_{1\times d_{\rm in}}
\underbrace{W}_{d_{\rm in}\times d_{\rm out}}
=\underbrace{\vect h^\top}_{1\times d_{\rm out}}.
\]
$W$ 的每一列规定一个输出怎样组合全部输入，这叫\textbf{线性映射}。
例如下一章用 $2\times3$ 的可训练矩阵，把两个输入数变成三个中间数。
文章表 $K$ 在本例中是给定数据；网络权重 $W$ 则由训练调整。

若遇到列向量写法，先转置即可对应：
$\vect h=W^\top\vect x$ 与 $\vect h^\top=\vect x^\top W$ 是同一次计算，
因为 $(AB)^\top=B^\top A^\top$。
主线继续使用行向量右乘权重。

\chapterchecktasks{03-matrix-multiplication}{理解与复算}{
\item \textbf{形状：}若 $A:2\times3$、$B:3\times4$，说明为什么 $AB$ 合法，并给出输出形状。若 $B$ 换为 $4\times3$ 呢？\chapterinlinehint{先核对内侧维度，再保留外侧维度。}
\item \textbf{手算：}沿用本章文章表 $K$，新需求为 $(1,0)$。求三个匹配分数，并说明为何要乘 $K^\top$。\chapterinlinehint{让需求行分别与三篇文章的特征列做点积。}
\item \textbf{找错：}某人断言“矩阵相乘和数字相乘一样，可以把 $AB$ 随意换成 $BA$”。用形状或一次反例指出问题。\chapterinlinehint{分别列出两个乘法的内侧维度。}
}

\par\smallskip
从一对向量的点积，到一个需求比较多篇文章，再到多个需求同时比较，
增加的只是结果的行数和列数。
下一章让右边的矩阵变成可训练权重，把这些运算接成一个多维网络。
